\chapter{Corpi non sferici}
Nel capitolo precedente abbiamo visto che un sistema caratterizzato da una densità $\rho(\vec{x}, t)$ soddisfa l'equazione di Poisson
\begin{align}\nabla^2 U(\vec{x}, t) = -4 \pi G \rho(\vec{x}, t)
\end{align}
che in questo capitolo ci prefiggiamo di risolvere nel caso in cui vi siano deviazioni dalla simmetria sferica. È istruttivo studiare inizialmente il caso con simmetria sferica perfetta per poi passare al caso più generale.
\section{Sistema a simmetria sferica}
Riscriviamo il laplaciano in coordinate sferiche e lo inseriamo nell'equazione di Poisson:
\begin{align}\left[ \frac{1}{r^2} \frac{\partial}{\partial r}\left(r^2 \frac{\partial}{\partial r}\right) + \frac{1}{r^2 \sin{\theta}} \frac{\partial}{\partial \theta}\left(\sin{\theta} \frac{\partial}{\partial \theta}\right) + \frac{1}{r^2 \sin^2{\theta}} \frac{\partial^2}{\partial \phi^2} \right] U = - 4 \pi G \rho.\end{align}
Se il sistema gode di simmetria sferica, il potenziale dipende esclusivamente dalla coordinata radiale $r$; pertanto l'equazione si riduce a
\begin{align}
    \frac{1}{r^2} \frac{\partial}{\partial r} \left( r^2 \frac{\partial U}{\partial r} \right) = -4 \pi G \rho(r, t),
\end{align}
da cui si ottiene
\begin{align}
    \frac{\partial U}{\partial r} = -\frac{4 \pi G}{r^2} \int_0^r \rho(r', t) (r')^2 dr' + \frac{C(t)}{r^2}.
\end{align}
La costante di integrazione $C(t)$ va posta a zero. Fisicamente, infatti, $\frac{\partial U}{\partial r}$ rappresenta il campo gravitazionale radiale. Nell'origine ($r=0$) la direzione radiale non è definita: una forza netta non nulla in quel punto punterebbe in una direzione preferenziale, rompendo la simmetria sferica del sistema. Inoltre, se $C(t)$ fosse diversa da zero, il termine divergerebbe. In definitiva si ha:
\begin{align}
    \frac{\partial U}{\partial r} = - \frac{4 \pi G m(r, t)}{r^2},
\end{align}
dove abbiamo definito la massa contenuta entro il raggio $r$:\begin{align}m(r, t) = \int_0^r \rho(r', t) (r')^2 dr'.\end{align}
Assumendo che il corpo sia confinato nello spazio, esiste un raggio $R$ tale per cui $\rho(r) = 0$ per ogni $r > R$. Pertanto, la massa totale
\begin{align}
    M \equiv m(R, t) = \int_0^R \rho(r', t) (r')^2 dr'
\end{align}
è una costante nel tempo. Lo si dimostra facilmente:
\begin{align}
    \frac{\mathrm{d} M}{\mathrm{d}t} = \int_0^R \frac{\partial \rho(r', t)}{\partial t} (r')^2 dr' = \int_0^R -\frac{1}{(r')^2} \nabla_r (\rho(r', t) v_r)(r')^2 dr' = 0,
\end{align}
dove abbiamo usato l'equazione di continuità in coordinate sferiche e il fatto che la densità e la velocità radiale $v_r$ siano nulle sul bordo del volume di integrazione. 
Abbiamo quindi mostrato che la massa totale del sistema si conserva. Inoltre, poiché per $r > R$ la densità si annulla, il campo gravitazionale all'esterno del corpo è dato da:
\begin{align}
    \frac{\partial U}{\partial r} = - \frac{4 \pi G M}{r^2} \implies U(r) = \frac{4 \pi G M}{r}.
\end{align}
Il potenziale all'esterno coincide con quello generato da una massa puntiforme $M$ posta nel centro del sistema: integrando l'equazione di Poisson abbiamo ottenuto una dimostrazione molto più diretta del teorema del guscio sferico di Newton.

\section{Deviazioni dalla simmetria sferica}
La simmetria sferica è una semplificazione utile, ma è raro che i corpi celesti siano perfettamente simmetrici. La rotazione attorno a un asse, l'interazione con altri corpi e le tensioni interne (come gli stress crostali delle stelle di neutroni) causano deviazioni da questa configurazione ideale. Quando tali deviazioni sono piccole, si procede tramite un'espansione in multipoli. Poiché il contributo dei momenti di ordine superiore tende a decrescere rapidamente, un troncamento dello sviluppo ai primi termini offre una buona descrizione del potenziale gravitazionale per la maggior parte dei problemi di interesse.
Nel caso della Terra, la misurazione dei momenti di multipolo fino a $l = 360$ (che quantificano le variazioni del campo gravitazionale su scale di centinaia di chilometri) ha permesso di studiare fenomeni come la deriva dei continenti, il vulcanismo e la struttura interna del pianeta. Questo approccio analitico, tuttavia, è limitato a sistemi approssimativamente sferici: le galassie a spirale, ad esempio, si discostano marcatamente da tale geometria, obbligando a ricorrere a metodi analitici differenti o a simulazioni numeriche per ricavarne il potenziale.
\subsection{Decomposizione in armoniche sferiche}
Ricordiamo che le armoniche sferiche $Y_{lm}$ sono autofunzioni dell'operatore differenziale angolare:
\begin{align}
    \left[ \frac{1}{\sin{\theta}} \frac{\partial}{\partial \theta} \left( \sin{\theta} \frac{\partial}{\partial \theta} \right) + \frac{1}{\sin^2{\theta}} \frac{\partial^2}{\partial \phi^2} \right] Y_{lm} = -l(l+1) Y_{lm},
\end{align}con $l$ intero non negativo e $m \in \{-l, \dots, l\}$ e riconosciamo, nel termine a sinistra, la parte angolare del laplaciano.
Dalla teoria di Sturm-Liouville sappiamo che queste funzioni costituiscono una base ortogonale e completa per lo spazio $L^2(S^2)$, ossia lo spazio delle funzioni a quadrato integrabile definite sulla sfera. Pertanto, vale la relazione di ortonormalità:
\begin{align}
    \int_{S^2} Y_{lm}^(\theta, \phi) Y_{l'm'}(\theta, \phi) \mathrm{d}\Omega = \delta_{ll'} \delta_{mm'},
\end{align}
e per ogni funzione $f \in L^2(S^2)$ si ha:
\begin{align}
f(\theta, \phi) = \sum_{l=0}^{\infty} \sum_{m=-l}^l f_{lm} Y_{lm}(\theta, \phi),
\end{align}
con i coefficienti dati da:
\begin{align}
f_{lm} = \int_{S^2} Y_{lm}(\theta, \phi) f(\theta, \phi) \mathrm{d}\Omega.
\end{align}
Possiamo quindi sviluppare sia la densità che il potenziale in serie di armoniche sferiche, ottenendo
\begin{align}
    &\rho(r, \theta, \phi, t) = \sum_{l=0}^{+\infty} \sum_{m=-l}^l \rho_{lm}(r, t) Y_{lm}(\theta, \phi), & &U(r, \theta, \phi, t) = \sum_{l=0}^{+\infty} \sum_{m=-l}^l U_{lm}(r, t) Y_{lm}(\theta, \phi),
\end{align}
dove i coefficienti dipendono da $r$ e $t$, poiché la proiezione viene effettuata su superfici sferiche a raggio e tempo fissati, e sono dati da
\begin{align}
    &\rho_{lm}(r, t) = \int_{S^2} Y_{lm}(\theta, \phi) \rho(r, \theta, \phi, t) \mathrm{d}\Omega, &
    &U_{lm}(r, t) = \int_{S^2} Y_{lm}(\theta, \phi) U(r, \theta, \phi, t) \mathrm{d}\Omega.
\end{align}

Inserendo questi sviluppi nell'equazione di Poisson si ottiene:
\begin{align}
    &\left[ \frac{1}{r^2} \frac{\partial}{\partial r}\left(r^2 \frac{\partial}{\partial r}\right) + \frac{1}{r^2 \sin{\theta}} \frac{\partial}{\partial \theta}\left(\sin{\theta} \frac{\partial}{\partial \theta}\right) + \frac{1}{r^2 \sin^2{\theta}} \frac{\partial^2}{\partial \phi^2} \right] \sum_{l=0}^{+\infty} \sum_{m=-l}^l U_{lm} Y_{lm} = \nonumber \\
    &= -4 \pi G \sum_{l=0}^{+\infty} \sum_{m=-l}^l \rho_{lm} Y_{lm}.\end{align}
Poiché le $Y_{lm}$ sono linearmente indipendenti, l'uguaglianza deve valere termine a termine per ogni coppia $(l, m)$, portando a:
\begin{align}
    \left[ \frac{\partial}{\partial r} \left( r^2 \frac{\partial}{\partial r} \right) - l(l+1) \right] U_{lm}(r, t) = - 4 \pi G r^2 \rho_{lm}(r, t).
\end{align}
In questo modo abbiamo ricondotto un'equazione alle derivate parziali a un sistema infinito di equazioni differenziali ordinarie disaccoppiate. Per risolvere queste equazioni utilizziamo il metodo delle funzioni di Green, definendo l'operatore differenziale
\begin{align}
    \mathcal{L} = \frac{\partial}{\partial r} r^2 \frac{\partial}{\partial r} - l(l+1),
\end{align}cerchiamo una funzione $g_l(r, r')$ che soddisfi
\begin{align}
    \mathcal{L} g_l(r, r') = - 4 \pi \delta(r-r').
\end{align}
Una volta determinata $g_l$, ogni coefficiente $U_{lm}$ potrà essere calcolato come
\begin{align} \label{eq:u_lm_green}
    U_{lm}(r, t) = G \int_0^\infty g_l(r, r') \rho_{lm}(r', t) (r')^2 dr'.
\end{align}
L'operatore $\mathcal{L}$ e, di conseguenza, la funzione di Green dipendono dall'indice $l$ ma non da $m$. Per $r \neq r'$, l'equazione per $g_l$ si riduce a quella omogenea:
\begin{align}
    \mathcal{L} g_l(r, r') = 0.
\end{align}
Cercando soluzioni della forma $r^\gamma$, si ha
\begin{align}
    \left[ \frac{\partial}{\partial r} r^2 \frac{\partial}{\partial r} - l(l+1) \right] r^\gamma = 0 \implies \gamma(\gamma+1) - l(l+1) = 0.
\end{align}
La combinazione lineare delle due radici indipendenti fornisce:\begin{align}g_l(r, r') = A(r')r^l + B(r')r^{-(l+1)}.\end{align}I coefficienti $A$ e $B$ dipendono da $r'$ poiché tale variabile gioca il ruolo di parametro per l'equazione differenziale in $r$. Affinché il potenziale sia fisicamente accettabile, non deve presentare divergenze spurie. Per $r < r'$, richiediamo che la soluzione sia regolare nell'origine ($r \to 0$); pertanto imponiamo $B = 0$:\begin{align}g_<(r, r') = A(r')r^l.\end{align}Analogamente, per $r > r'$, richiediamo che il potenziale si annulli all'infinito ($r \to \infty$); pertanto imponiamo $A = 0$:
\begin{align}
    g_>(r, r') = B(r')r^{-(l+1)}.
\end{align}
Combinando i risultati mediante il gradino di Heaviside, si ottiene:
\begin{align}
    g_l(r, r') = A(r')r^l \theta(r'-r) + B(r')r^{-(l+1)}\theta(r-r').
\end{align}
Imponendo infine la continuità della funzione di Green in $r = r'$, si ricava la condizione
\begin{align}
    A(r')(r')^l = B(r')(r')^{-(l+1)},
\end{align}
dunque possiamo riscrivere nella seguente forma
\begin{align}
    g_l(r, r') = C(r') \left[ \frac{r^l}{(r')^l}\theta(r'-r) + \frac{r^{-(l+1)}}{(r')^{-(l+1)}}\theta(r-r') \right].
\end{align}
ed integrando l'equazione differenziale fra $[r'-\varepsilon, r'+\varepsilon]$ con $\varepsilon \to 0$ si ottiene che
\begin{align}
    \int_{r'-\varepsilon}^{r'+\varepsilon} \left[ \frac{\partial}{\partial r} r^2 \frac{\partial}{\partial r} \right] g_l(r, r') = - 4 \pi.
\end{align} 
Si noti che il termine proporzionale all'integrale di $g_l(r, r')$ si annulla per continuità nel limite di $\varepsilon \to 0$: rimane dunque
\begin{align}
    \int_{r'-\varepsilon}^{r'+\varepsilon} \left[ \frac{\partial}{\partial r} \left( r^2 \frac{\partial g_l(r, r')}{\partial r} \right) \right] = - 4 \pi,
\end{align}
in cui si ottiene che
\begin{align}
    C(r')[(-l+1)(r') - (l+2)(r')] = -4 \pi \implies C(r') = \frac{4 \pi}{2l+1} \frac{1}{r'}.
\end{align}
La funzione di Green dell'operatore $\mathcal{L}$ risulta, dunque, essere
\begin{align}
    g_l(r, r') = \frac{4 \pi}{2l+1} \left[ \frac{r^l}{(r')^{l+1}} \theta(r' - r) + \frac{(r')^{l}}{r^{l+1}} \theta(r-r') \right],
\end{align}
che può essere riscritta in forma più compatta introducendo $r_> = \max\{r, r'\}$ e $r_< = \min\{r, r'\}$, ottenendo
\begin{align}
    g_l(r, r') = \frac{4 \pi}{2l+1} \frac{r_<^l}{r_>^{l+1}}.
\end{align}
Sostituendo nella \eqref{eq:u_lm_green} si ottiene che
\begin{align}
    U_{lm}(r, t) &= \frac{4 \pi G}{2l+1} \int_{0}^{+\infty} \frac{r^l_<}{r_>^{l+1}} \rho_{lm}(r', t) \mathrm{d}r' = \nonumber \\
    &= \frac{4 \pi G}{2l+1} \left[ r^l \int_r^{+\infty} \frac{\rho_{lm}(r', t)}{(r')^{l+1}} (r')^2 dr' + \frac{1}{r^{l+1}} \int_0^r (r')^l \rho_{lm}(r', t)(r')^2 \mathrm{d}r' \right].
\end{align}

